(1) 在原点外，球对称方程 \((r^2\phi')'-\kappa^2r^2\phi=0\)，令 \(u=r\phi\) 得 \(u''-\kappa^2u=0\)。无穷远衰减只留 \(e^{-\kappa r}\)；原点小球高斯通量 \(-4\pi\varepsilon r^2\phi'\to q\) 定系数，故 \(\phi(r)=q e^{-\kappa r}/(4\pi\varepsilon r)\)，\(\boldsymbol E=-\nabla\phi=[q/(4\pi\varepsilon r^2)](1+\kappa r)e^{-\kappa r}\hat{\boldsymbol r}\)。

(2) 感应密度 \(\rho_{\rm ind}(r)=-q\kappa^2e^{-\kappa r}/(4\pi r)\)，在原点附近虽按 \(1/r\) 发散，但体积积分可积。半径 \(r\) 内感应电荷为 \(Q_{\rm ind}(r)=-q\kappa^2\int_0^r s e^{-\kappa s}ds=-q[1-(1+\kappa r)e^{-\kappa r}]\)。因此 \(Q_{\rm ind}(\infty)=-q\)，屏蔽云恰好抵消测试电荷的总电量。

(3) 当 \(r\gg\kappa^{-1}\)，电势具有 \(e^{-\kappa r}/r\) 的 Yukawa 尾而非无屏蔽的 \(1/r\) 尾。当固定 \(r>0\) 令 \(\kappa\to0^+\) 时，\(\phi\to q/(4\pi\varepsilon r)\)，\(Q_{\rm ind}(r)\to0\)。需要区分在固定半径下撤去屏蔽与先取无限远：两极限一般不可对调，因为总屏蔽云被推向无限远。

(4) 写 \(\phi=\phi_{\rm bare}+\phi_{\rm ind}\)，其中 \(\phi_{\rm bare}=q/(4\pi\varepsilon r)\)。在原点取有限诱导部分 \(\phi_{\rm ind}(0)=\lim_{r\to0}{q\over4\pi\varepsilon r}(e^{-\kappa r}-1)=-q\kappa/(4\pi\varepsilon)\)。线性响应充电过程的耦合自由能为 \(\Delta F=\tfrac12 q\phi_{\rm ind}(0)=-q^2\kappa/(8\pi\varepsilon)\)。直接用 \(\tfrac12q\phi(0)\) 会因裸点电荷自能发散而失去物理意义。
